\section{按尺度截断的误差}
\label{sec:besov-scale-truncation}

小波坐标中最直接的近似方法，是对每个函数保留同一个尺度范围内的全部系数，舍去更细的层。本节固定$p=2$，以$L^2$范数度量误差，先写出截断误差的精确公式，再对Besov球给出不依赖函数的统一截断界。最后反过来说明，各个截断尺度上的误差汇总在一起，恰好刻画Besov条件，这也解释了参数$q$在逼近问题中的意义。

$p=2$时，式\eqref{eq:besov-wavelet-norm}的权重化为$2^{js}$，于是由定理~\ref{thm:besov-wavelet-characterization}，
\begin{equation}
C_0^{-1}\Bigl(\lVert a\rVert_{\ell^2}+\bigl\lVert(2^{js}\lVert c_j\rVert_{\ell^2})_{j\geq j_0}\bigr\rVert_{\ell^q}\Bigr)\leq\lVert f\rVert_{B^s_{2,q}}\leq C_0\Bigl(\lVert a\rVert_{\ell^2}+\bigl\lVert(2^{js}\lVert c_j\rVert_{\ell^2})_{j\geq j_0}\bigr\rVert_{\ell^q}\Bigr).
\label{eq:besov-wavelet-norm-l2}
\end{equation}
本节与后面两节都使用这一形式。

\subsection{尺度截断}
\label{sec:besov-truncation}

对整数$J\geq j_0$，令$V_J$为粗项基函数与$j_0\leq j<J$的全部细节小波张成的子空间，$P_J$为到$V_J$的正交投影：
\[
P_Jf=\sum_{k}a_k\varphi_{j_0,k}+\sum_{j_0\leq j<J}\sum_kc_{j,k}\psi_{j,k} .
\]
$V_J$的维数为$M_J=2^{j_0}+\sum_{j=j_0}^{J-1}2^j=2^J$。由正交性与Parseval等式，截断误差
\begin{equation}
E_J(f)=\lVert f-P_Jf\rVert_2,\qquad E_J(f)^2=\sum_{j\geq J}\lVert c_j\rVert_{\ell^2}^2 .
\label{eq:besov-truncation-error}
\end{equation}
由第~\ref{chap:inner-product-spaces}章，$P_Jf$是$V_J$中离$f$最近的元素，所以$E_J(f)$也是在$V_J$中能达到的最佳误差。

这里有两个不同的量需要分开：尺度$J$规定保留到多细的层，项数预算$M_J$规定用了多少个基函数。二者由$M_J=2^J$联系起来，但含义不同，下一节的最佳$N$项逼近只限制项数而不限制尺度。更重要的是，$V_J$在看到任何函数之前就已经选定，对整类函数都相同。这正是第~\ref{chap:function-approximation}章所说的固定子空间逼近：近似是线性的，所有函数共用同一份基函数表。

\subsection{统一截断误差}
\label{sec:besov-uniform-truncation}

现在对Besov球中的所有函数同时估计截断误差。

\begin{theorem}[Besov球上的统一截断界]
\label{thm:besov-uniform-truncation}
设$s>0$，$1\leq q\leq\infty$。存在只依赖$s$与$C_0$的常数$C_s$，使对每个满足$\lVert f\rVert_{B^s_{2,q}}\leq R$的$f$和每个$J\geq j_0$，
\[
E_J(f)\leq C_sR\,2^{-Js}=C_sR\,M_J^{-s} .
\]
\end{theorem}

\begin{proof}
由式\eqref{eq:besov-wavelet-norm-l2}，$\ell^q$范数不小于各项的最大值，所以对每个$j\geq j_0$，$2^{js}\lVert c_j\rVert_{\ell^2}\leq C_0R$，即$\lVert c_j\rVert_{\ell^2}\leq C_0R2^{-js}$。代入式\eqref{eq:besov-truncation-error}，对几何级数求和，
\[
E_J(f)^2\leq C_0^2R^2\sum_{j\geq J}2^{-2js}=\frac{C_0^2R^2\,2^{-2Js}}{1-2^{-2s}} .
\]
取$C_s=C_0(1-2^{-2s})^{-1/2}$，再用$2^J=M_J$即得。
\end{proof}

常数$C_s$不随函数和截断尺度改变，所以这是整类函数的最坏误差界：用$M$个预先固定的基函数，Besov球中每个函数的误差都不超过$C_sRM^{-s}$。这个量级对整类函数不能再改进。取$f=(R/C_0)2^{-Js}\psi_{J,0}$，它的系数范数为$R/C_0$，所以$\lVert f\rVert_{B^s_{2,q}}\leq R$；而它的全部能量都在第$J$层，$E_J(f)=(R/C_0)2^{-Js}$，与上界只差常数倍。注意这个界只用到每层的统一上界，与$q$无关：$q$越大球越大，但每层的最大可能大小相同。

\subsection{截断误差的尺度汇总}
\label{sec:besov-truncation-aggregation}

统一截断界是从Besov条件推出误差。反过来，一个函数在各个截断尺度上的误差，能否读出它属于哪个Besov空间？答案是肯定的，但要把各尺度的误差汇总起来，而不是只看一个速率。

\begin{theorem}[截断误差刻画Besov条件]
\label{thm:besov-truncation-characterization}
设$s>0$，$1\leq q\leq\infty$。存在常数$C$，使对每个$f\in L^2(\mathbb T)$，
\[
C^{-1}\lVert f\rVert_{B^s_{2,q}}\leq\lVert f\rVert_2+\bigl\lVert(2^{Js}E_J(f))_{J\geq j_0}\bigr\rVert_{\ell^q}\leq C\lVert f\rVert_{B^s_{2,q}} ,
\]
两端同时有限或同时无限。
\end{theorem}

\begin{proof}
记$b_j=2^{js}\lVert c_j\rVert_{\ell^2}$，$e_J=2^{Js}E_J(f)$。一个方向来自误差至少包含当前层：由式\eqref{eq:besov-truncation-error}，$\lVert c_J\rVert_{\ell^2}\leq E_J(f)$，所以$b_J\leq e_J$对每个$J$成立。

另一方向来自尾和。平方和的平方根不超过和，所以$E_J(f)\leq\sum_{j\geq J}\lVert c_j\rVert_{\ell^2}$，乘以$2^{Js}$得
\[
e_J\leq\sum_{m\geq0}2^{-ms}b_{J+m} .
\]
对序列$(e_J)_J$取$\ell^q$范数，由Minkowski不等式，右端每个$m$的贡献不超过$2^{-ms}\lVert(b_j)\rVert_{\ell^q}$，于是$\lVert(e_J)\rVert_{\ell^q}\leq(1-2^{-s})^{-1}\lVert(b_j)\rVert_{\ell^q}$。

最后处理粗项：$\lVert a\rVert_{\ell^2}\leq\lVert f\rVert_2$，并且$\lVert f\rVert_2\leq\lVert a\rVert_{\ell^2}+E_{j_0}(f)\leq\lVert a\rVert_{\ell^2}+2^{-j_0s}e_{j_0}$。结合式\eqref{eq:besov-wavelet-norm-l2}即得两侧估计。
\end{proof}

所以$B^s_{2,q}$中的$q$在逼近语言中有直接含义：它规定截断误差序列$(2^{Js}E_J(f))_J$怎样汇总。$q=\infty$要求统一速率$E_J(f)\leq C2^{-Js}$，有限的$q$还要求这些按速率放大的误差跨尺度$q$次可和。

\begin{example}[统一速率与跨尺度可和性]
\label{ex:besov-one-coefficient-per-scale}
令$f=\sum_{j\geq j_0}2^{-js}\psi_{j,0}$，即每层只有一个大小为$2^{-js}$的系数。系数平方可和，$f\in L^2(\mathbb T)$。这时$b_j=1$对每个$j$成立，序列$(b_j)$有界而不$q$次可和；由式\eqref{eq:besov-wavelet-norm-l2}，$f\in B^s_{2,\infty}$，但对任何有限$q$，$f\notin B^s_{2,q}$。截断误差为
\[
E_J(f)=\Bigl(\sum_{j\geq J}2^{-2js}\Bigr)^{1/2}=\frac{2^{-Js}}{(1-2^{-2s})^{1/2}},
\]
恰好具有统一速率$2^{-Js}$。所以仅凭统一衰减速率，只能保证$q=\infty$的条件；有限$q$所需的跨尺度可和性，需要误差在越来越细的尺度上比$2^{-Js}$衰减得更快，并且快到可以$q$次求和。这与第~\ref{sec:besov-distribution-aggregation}小节$b_j=1$与$b_j=1/j$的对比是同一件事。
\end{example}
